\section{Análisis horizontal}

\begin{frame}{Análisis horizontal: definición}
  \begin{columns}[T,onlytextwidth]
    \column{0.45\textwidth}
    \footnotesize
    El análisis horizontal compara una misma partida contable entre dos o más periodos.

    \vspace{0.5em}
    Permite responder cuánto cambió una cuenta respecto al periodo base y si ese cambio es favorable o requiere revisión.
    \tiempo{0.7 minutos}
    \column{0.51\textwidth}
    \figpng[0.66\textheight]{F04_Analisis_Horizontal_Linea_Tiempo.png}
  \end{columns}
  \note{Definir análisis horizontal.}
\end{frame}

\begin{frame}{Utilidad del análisis horizontal}
  \begin{columns}[T,onlytextwidth]
    \column{0.45\textwidth}
    \begin{itemize}
      \item Detectar crecimiento.
      \item Reconocer disminuciones.
      \item Identificar tendencias.
      \item Evaluar señales de alerta.
    \end{itemize}
    \micro{Su valor está en explicar la evolución de las partidas y no solo en mostrar que aumentaron o disminuyeron.}
    \tiempo{0.7 minutos}
    \column{0.51\textwidth}
    \figpng[0.62\textheight]{F07_Interpretacion_Horizontal.png}
  \end{columns}
  \note{Explicar utilidad del horizontal.}
\end{frame}

\begin{frame}{Fórmulas del análisis horizontal}
  \begin{columns}[T,onlytextwidth]
    \column{0.46\textwidth}
    \[
      VA = V_{actual} - V_{base}
    \]
    \[
      VP = \frac{VA}{V_{base}} \times 100
    \]
    \begin{itemize}
      \item VA: variación absoluta.
      \item VP: variación porcentual.
    \end{itemize}
    \tiempo{0.8 minutos}
    \column{0.50\textwidth}
    \figpng[0.64\textheight]{F05_Formula_Analisis_Horizontal.png}
  \end{columns}
  \note{Explicar fórmulas.}
\end{frame}

\begin{frame}{Ejemplo básico de cálculo horizontal}
  \begin{columns}[T,onlytextwidth]
    \column{0.48\textwidth}
    \footnotesize
    Una empresa desea comparar sus ventas entre dos periodos para medir el crecimiento comercial.
    \vspace{0.35em}
    \begin{tabular}{lr}
      \toprule
      Concepto & Monto \\
      \midrule
      Ventas 2024 & S/ 100,000.00 \\
      Ventas 2025 & S/ 120,000.00 \\
      Variación absoluta & S/ 20,000.00 \\
      Variación porcentual & 20.00 \% \\
      \bottomrule
    \end{tabular}
    \micro{Las ventas aumentaron S/ 20,000.00, equivalentes al 20.00 \% del periodo base.}
    \tiempo{0.8 minutos}
    \column{0.48\textwidth}
    \figpng[0.58\textheight]{F05_Formula_Analisis_Horizontal.png}
  \end{columns}
  \note{Explicar ejemplo básico.}
\end{frame}

\begin{frame}{Ejemplo aplicado: variación de costo cloud}
  \begin{columns}[T,onlytextwidth]
    \column{0.42\textwidth}
    \footnotesize
    Una empresa de servicios digitales contrata infraestructura cloud para alojar su plataforma.
    \vspace{0.25em}
    \begin{itemize}
      \item Costo 2024: S/ 8,000.00.
      \item Costo 2025: S/ 12,000.00.
      \item Variación: S/ 4,000.00 y 50.00 \%.
    \end{itemize}
    \micro{Si los ingresos no crecieron en la misma proporción, debe revisarse arquitectura, proveedor o plan contratado.}
    \tiempo{0.8 minutos}
    \column{0.54\textwidth}
    \figpng[0.66\textheight]{F18_Ejemplo_Cloud_Horizontal.png}
  \end{columns}
  \note{Ejemplo cloud.}
\end{frame}

\begin{frame}{Procedimiento del análisis horizontal}
  \begin{columns}[T,onlytextwidth]
    \column{0.46\textwidth}
    \begin{enumerate}
      \item Seleccionar partida.
      \item Elegir periodo base.
      \item Calcular variación absoluta.
      \item Calcular variación porcentual.
      \item Interpretar resultado.
    \end{enumerate}
    \tiempo{0.8 minutos}
    \column{0.50\textwidth}
    \figpng[0.62\textheight]{F04_Analisis_Horizontal_Linea_Tiempo.png}
  \end{columns}
  \note{Explicar procedimiento horizontal.}
\end{frame}

\begin{frame}{Caso Empresa ABC: datos para análisis horizontal}
  \begin{columns}[T,onlytextwidth]
    \column{0.52\textwidth}
    \scriptsize
    \begin{tabular}{lrr}
      \toprule
      Partida & 2024 & 2025 \\
      \midrule
      Ventas netas & S/ 180,000.00 & S/ 215,000.00 \\
      Costo ventas & S/ 108,000.00 & S/ 133,000.00 \\
      Utilidad bruta & S/ 72,000.00 & S/ 82,000.00 \\
      Utilidad neta & S/ 21,855.00 & S/ 25,027.00 \\
      \bottomrule
    \end{tabular}
    \micro{Estas partidas permiten observar crecimiento comercial, comportamiento del costo y evolución de utilidad.}
    \tiempo{0.8 minutos}
    \column{0.44\textwidth}
    \figpng[0.58\textheight]{F06_Barras_Comparativas_2024_2025.png}
  \end{columns}
  \note{Presentar datos ABC.}
\end{frame}

\begin{frame}{Caso Empresa ABC: cálculo horizontal}
  \begin{columns}[T,onlytextwidth]
    \column{0.48\textwidth}
    \scriptsize
    \begin{tabular}{lrrr}
      \toprule
      Partida & 2024 & 2025 & Var. \% \\
      \midrule
      Ventas & S/ 180,000.00 & S/ 215,000.00 & 19.44 \% \\
      Costo & S/ 108,000.00 & S/ 133,000.00 & 23.15 \% \\
      U. bruta & S/ 72,000.00 & S/ 82,000.00 & 13.89 \% \\
      U. neta & S/ 21,855.00 & S/ 25,027.00 & 14.51 \% \\
      \bottomrule
    \end{tabular}
    \tiempo{0.8 minutos}
    \column{0.48\textwidth}
    \figpng[0.62\textheight]{F06_Barras_Comparativas_2024_2025.png}
  \end{columns}
  \note{Explicar cálculo horizontal del caso.}
\end{frame}

\begin{frame}{Caso Empresa ABC: interpretación horizontal}
  \begin{columns}[T,onlytextwidth]
    \column{0.46\textwidth}
    \footnotesize
    Las ventas aumentan 19.44 \%, pero el costo de ventas aumenta 23.15 \%.

    \vspace{0.5em}
    Esto indica que el crecimiento comercial no se convierte totalmente en eficiencia operativa, porque el costo avanza con mayor intensidad.
    \tiempo{0.8 minutos}
    \column{0.50\textwidth}
    \figpng[0.64\textheight]{F07_Interpretacion_Horizontal.png}
  \end{columns}
  \note{Interpretar caso horizontal.}
\end{frame}

\begin{frame}{Decisión empresarial derivada del análisis horizontal}
  \begin{columns}[T,onlytextwidth]
    \column{0.48\textwidth}
    \begin{itemize}
      \item Revisar proveedores y costos.
      \item Evaluar precios y productividad.
      \item Mejorar procesos operativos.
      \item Configurar alertas en dashboards.
    \end{itemize}
    \micro{Si el costo crece más rápido que las ventas, la gerencia debe actuar antes de que se deteriore la rentabilidad.}
    \tiempo{0.7 minutos}
    \column{0.48\textwidth}
    \figpng[0.62\textheight]{F14_Decisiones_Empresariales.png}
  \end{columns}
  \note{Decisión derivada horizontal.}
\end{frame}
